% English version of paper/ja_v4/02_関連研究.md (ver4, 2026-11-17): three paragraphs, each closing with the difference from this study. Cited claims carried over from v3_02_related.tex; the regrouped sentences follow reviews/v4_ja/citations_2026-11-17.md.
\section{Related work}\label{sec:related}

\subsection{Adapting classifiers to a target with few labels}\label{sec:rel-da}
Learning representations that do not reveal the domain (domain-adversarial neural networks, DANN) \cite{ganin2016dann} and aligning feature covariances (CORAL) \cite{sun2016coral} have been proposed for unlabelled targets, entropy minimisation and pseudo-labels for semi-supervised learning without domain shift \cite{grandvalet2004,lee2013pseudolabel}, and minimax entropy (MME) for semi-supervised domain adaptation, in which the target has a few labels \cite{saito2019mme}. Adding source data can also lower target performance (negative transfer) \cite{wang2019negative}. In emotion recognition, adversarial learning that aligns representations with unlabelled target utterances \cite{abdelwahab2018dann} and learning that pulls together representations learned per corpus \cite{gideon2021addog} have been proposed for audio classifiers; for targets with a few labels, Feng and Chaspari \cite{feng2023} proposed metric learning for audio features, Li et al.\ \cite{li2021acted} adaptation of audio and video between acted and improvised speech, and He et al.\ \cite{he2025m4ser} a method using audio and speech-recognition transcripts. These domain adaptation and semi-supervised methods, and the adaptation methods for emotion recognition, all adapt the classifier with data of the same modality as the classifier being retrained. MER fusion methods, including those for sentiment polarity, have instead been evaluated by training and testing within one corpus \cite{zadeh2017tfn,mittal2020m3er,yu2021selfmm,yun2024telme}, and the cross-corpus evaluations of Ryumina et al.\ \cite{ryumina2025} and Lian et al.\ \cite{lian2026merbench} did not use target labels. We compare 11 methods, including DANN, CORAL, entropy minimisation, self-training and MME (Sect.~\ref{sec:distill}), with pseudo-labels including other modalities under the same conditions.

\subsection{Learning from predictions of other modalities}\label{sec:rel-distill}
Albanie et al.\ \cite{albanie2018} attached facial-expression predictions to unlabelled videos and trained an audio emotion classifier on them, and Zhang et al.\ \cite{zhang2021combining} combined training towards targets that include predictions of other modalities with semi-supervised learning. Transferring representations or predictions of one modality to a model of another through paired unlabelled data is also used outside emotion recognition \cite{gupta2016,jaritz2020xmuda}. Training on the predictions of a more accurate model of another modality has been reported not to always improve the trained model \cite{xue2023mfh}. Three comparisons are close to ours. Zhang et al.\ (2021) compared pseudo-labels including facial predictions with same-modality semi-supervised learning for speech emotion recognition, but used all training labels of the evaluated corpus (IEMOCAP) and did not compare domain adaptation methods. Jaritz et al.\ \cite{jaritz2023xmossda} compared cross-modal learning with uni-modal unsupervised domain adaptation methods moved to a semi-supervised setting, but for 3D semantic segmentation. Zhang et al.\ (2024) \cite{zhang2024amanda} proposed a multimodal domain adaptation method that uses no target labels between IEMOCAP, MELD, MOSEI and MSP-IMPROV and compared it with DANN and other methods, but not with distillation on pseudo-labels from other modalities. To the best of our knowledge, no prior study has compared, in cross-corpus MER with audio, visual and text and only a few target labels per emotion, distillation on pseudo-labels including other modalities with 11 domain adaptation, semi-supervised and test-time adaptation methods that do not use such pseudo-labels (versions moved to linear classifiers on frozen features) under the same conditions.

\subsection{Late fusion and reliability weights}\label{sec:rel-fusion}
Zadeh et al.\ \cite{zadeh2017tfn} criticised late fusion, which trains unimodal classifiers separately and combines their predictions, for not learning interactions between modalities because it assumes that simple weighted averaging is a proper fusion approach. Predictions can be combined by product or sum rules \cite{kittler1998}. Busso et al.\ \cite{busso2004} reported that modality strengths differ by emotion. Kaya et al.\ \cite{kaya2017} chose model-by-emotion weights on AFEW by random search on validation data, M3ER \cite{mittal2020m3er} learned a multiplicative fusion layer that strengthens reliable modality cues per utterance, and Mar\'in et al.\ \cite{marin2026classwise} used per-utterance, per-emotion modality weights output by a network trained with all labels within one corpus. Kaya et al.\ determined the weights, and M3ER and Mar\'in et al.\ the layers that output them, with labelled data of the evaluated corpus. With few training samples, a model trained on more modalities has also been argued to be able to fall below a model trained on a subset of them \cite{huang2021}. Like Kaya et al., the proposal uses one weight per classifier and emotion, but it sets the weights from a few target labels and the source corpora. We found no study that splits the weights into classifier strengths measured on a few target labels per emotion and per-emotion profiles measured on the source corpora.
